% L'ellipse de la hessienne en un pixel d'une crête sombre.
% Axes portés par les vecteurs propres e_1, e_2 ; demi-axes 1/sqrt(|lambda_k|).
% Crête : |lambda_1| petit (grand axe, le long de la fissure, angle theta),
% lambda_2 > 0 grand (petit axe, en travers).
% La boîte englobante est tenue serrée : le rendu passe par un \resizebox.
    \begin{tikzpicture}[x=1cm,y=1cm,scale=1.0,every node/.style={scale=1.0},
        crete/.style={line width=9pt, gris_fonce_inria!55, line cap=round},
        vecp/.style={-{Latex[length=2.4mm]}, line width=1.1pt},
        filet/.style={line width=0.8pt}]

      \begin{scope}[rotate=27]
        \draw[crete] (-2.05,0)--(2.05,0);
        \draw[inria-rouge, line width=1pt, fill=inria-rouge!8]
              (0,0) ellipse (1.75 and 0.62);
        % longueurs corrigées de la pointe : les pointes tombent en 1,75 et 0,62
        \draw[vecp, inria-rouge]    (0,0)--(1.833,0);
        \draw[vecp, inria-rouge!62] (0,0)--(0,0.703);
        \draw[inria-rouge!40, dashed, line width=0.6pt] (1.75,0)--(2.30,0);
        \node[font=\small, inria-rouge, fill=white, inner sep=1pt] at (0.95,0) {$e_1$};
        \node[font=\small, inria-rouge!62, fill=white, inner sep=1pt] at (0,0.40) {$e_2$};
      \end{scope}

      % axe de référence et angle theta, tracés hors de l'ellipse
      \draw[gray!65, dashed, line width=0.6pt] (0,0)--(2.20,0);
      \draw[gris_fonce_inria!80, line width=0.7pt] (1.98,0) arc (0:27:1.98);
      \node[font=\small] at (1.82,0.45) {$\theta$};

      \node[circle,fill=inria-2024-bleu-canard,inner sep=0pt,minimum size=4.2pt] (px) at (0,0) {};
      \node[font=\small, below left=0pt and 1pt of px] {$x$};

      % demi-axes
      \draw[inria-rouge, filet] (1.5591,0.7946)--(1.66,1.04);
      \node[font=\small, inria-rouge, anchor=south] at (1.66,1.03)
            {$\dfrac{1}{\sqrt{|\lambda_1|}}$};
      \draw[inria-rouge!62, filet] (-0.2815,0.5524)--(-0.62,0.90);
      \node[font=\small, inria-rouge!62, anchor=south east] at (-0.58,0.88)
            {$\dfrac{1}{\sqrt{|\lambda_2|}}$};
    \end{tikzpicture}
